\documentclass[a4paper,12pt,oneside]{book}

%-------------------------------Start of the Preable------------------------------------------------
\usepackage[english]{babel}
\usepackage{blindtext}
%packagr for hyperlinks
\usepackage{hyperref}
\hypersetup{
    colorlinks=true,
    linkcolor=blue,
    filecolor=magenta,      
    urlcolor=cyan,
}

\urlstyle{same}
%use of package fancy header
\usepackage{fancyhdr}
\setlength\headheight{26pt}
\fancyhf{}
%\rhead{\includegraphics[width=1cm]{logo}}
\lhead{\rightmark}
\rhead{\includegraphics[width=1cm]{logo.png}}
\fancyfoot[RE, RO]{\thepage}
\fancyfoot[CE, CO]{\href{http://www.e-yantra.org}{www.e-yantra.org}}

\pagestyle{fancy}

%use of package for section title formatting
\usepackage{titlesec}
\titleformat{\chapter}
  {\Large\bfseries} % format
  {}                % label
  {0pt}             % sep
  {\huge}           % before-code
 
%use of package tcolorbox for colorful textbox
\usepackage[most]{tcolorbox}
\tcbset{colback=cyan!5!white,colframe=cyan!75!black,halign title = flush center}

\newtcolorbox{mybox}[1]{colback=cyan!5!white,
colframe=cyan!75!black,fonttitle=\bfseries,
title=\textbf{\Large{#1}}}

%use of package marginnote for notes in margin
\usepackage{marginnote}

%use of packgage watermark for pages
%\usepackage{draftwatermark}
%\SetWatermarkText{\includegraphics{logo}}
\usepackage[scale=2,opacity=0.1,angle=0]{background}
\backgroundsetup{
contents={\includegraphics{logo}}
}

%use of newcommand for keywords color
\usepackage{xcolor}
\newcommand{\keyword}[1]{\textcolor{red}{\textbf{#1}}}

%package for inserting pictures
\usepackage{graphicx}

%package for highlighting
\usepackage{color,soul}

%new command for table
\newcommand{\head}[1]{\textnormal{\textbf{#1}}}


%----------------------End of the Preamble---------------------------------------


\begin{document}

%---------------------Title Page------------------------------------------------
\begin{titlepage}
\raggedright
{\Large eYSIP2017\\[1cm]}
{\Huge\scshape Vegetable Identification Using Transfer Learning \\[.1in]}
\vfill
\begin{flushright}
{\large Sanket Shanbhag \\}
{\large Supriya Suresh \\}
{\large Saurav and Naveen \\}
{\large Duration of Internship: $ 22/05/2017-07/07/2017 $ \\}
\end{flushright}

{\itshape 2017, e-Yantra Publication}
\end{titlepage}
%-------------------------------------------------------------------------------

\chapter[Project Tag]{Vegetable Identification Using \\ Transfer Learning}
\section*{Abstract}
This project aims to create a system for automatically logging farm produce from the green house at e-Yantra. An image is captured of the produce and this image is fed to a neural network based on the Inception-v3 model by Google. This system then identifies the vegetable and classifies it into different classes it has been trained on. Also, a re-training system is implemented which automatically trains any new images captured and thereby further increasing the accuracy of the system over time.

\subsection*{Completion status}
The system has been successfully trained for 18 classes (vegetables) achieving accuracy of 99.5\% on a separate test set. Also a system has been implemented for auto-training any new images captured on the system every 4 days.

\section{Hardware parts}
\begin{itemize}
  \item List of hardware 
  \item Detail of each hardware: \href[page=5]{./datasheet/MPU-9150.pdf}{Datasheet, page 5}, \href{http://www.amazon.in}{Vendor link}, 
  \item Connection diagram
\end{itemize}

\section{Software used}
\begin{itemize}
  \item \href{https://www.python.org/downloads/release/python-352/}{Python 3.5.2}
  \item \href{https://pypi.python.org/pypi/tensorflow/1.2.0}{Tensorflow 1.2.0}
  \item \href{https://github.com/google/prettytensor}{PrettyTensor 0.7.4} 
  \item Installation Steps:
  \begin{itemize}
  	\item Open a terminal in the project base folder and type: \\
  	\texttt{pip3 install -r requirements.txt}
  \end{itemize}
\end{itemize}

\section{Assembly of hardware}
Circuit diagram and Steps of assembly of hardware with pictures for each step
\subsection*{Circuit Diagram}
Circuit schematic, simplified circuit diagram , block diagram of system
\subsection*{Step 1}
Steps for assembling part 1
\subsection*{Step 2}
Steps for assembling part 2
\subsection*{Step 3}
Steps for assembling part 3



\section{Software and Code}
The complete code is available \href{https://github.com/eYSIP-2017/eYSIP-2017_Vegetable-Identification-Using-Transfer-Learning}{here}. It is divided into various folders each with a self contained module of the project.

\subsection{Downloading and formatting data}
The folder \texttt{download\_data} contains python scripts and shell scripts for downloading and formatting data from a list of URLS or from \href{http://www.image-net.org/}{ImageNet}. Information on using these scripts can be found in the \href{https://github.com/eYSIP-2017/eYSIP-2017_Vegetable-Identification-Using-Transfer-Learning/wiki/Downloading-and-formatting-data.}{wiki page}.

\subsection{Transfer Learning}
The code for transfer learning can be found in the \texttt{transfer\_on\_inception\_v3} folder. The \texttt{transferveg.py} file contains the code for running the model.
We remove the final layer of the Inception-v3 model and add 3 fully connected layers of 4096, 2048, and 1024 nodes with a dropout layer after the 2048-node layer.
This helps us achieve an accuracy of 99.5\%

\subsection{Module Integration}
The \texttt{ghfarm.py} script contained in the \texttt{module\_integration} folder will run on the raspberry pi on the weighing machine. It interfaces with the server to predict the crop and sends data once the image is taken.

\subsection{Auto-Training}
The \texttt{AutoTrain} folder contains code for running a server to accept images and scripts for adding cron tasks to run the autotraining code every 4 days.


\section{Use and Demo}
Final Setup Image

User Instruction for demonstration

\href{http://www.youtube.com}{Youtube Link} of demonstration video 

\section{Future Work}
What can be done to take this work ahead in future as projects.

\section{Bug report and Challenges}
Any issues in code and hardware.

Any failure or challenges faced during project

\begin{thebibliography}{li}
\bibitem{wavelan97}
Ad Kamerman and Leo Monteban,
{\em WaveLAN-II: A High-Performance Wireless LAN for the Unlicensed band},
1997.

\end{thebibliography}


\end{document}

